\documentclass{article}
% Requires paper-utils-full.sty, copied here as paper-utils.sty.
% To hide notes, uncomment before loading paper-utils:
% \PassOptionsToPackage{disable}{todonotes}
\usepackage{paper-utils}
\setlength{\marginparwidth}{2cm}

\title{A Brief Guide to the \texttt{todonotes} package}
\date{\today}
\author{\href{https://pages.di.unipi.it/massa}{Jacopo Massa}}

\begin{document}

\maketitle

\section*{\texttt{todo} command}

Use the \verb|\todo{<text>}| command to insert a note in the margin. The note will be numbered and will appear in the list of todos at the end of the document.
\begin{Verbatim}
	This paragraph needs a reference.\todo{Add a reference.}
\end{Verbatim}
Output:
This paragraph needs a reference.\todo{Add a reference.}

\section*{Inline notes}

Use the \verb|\todo[inline]{<text>}| command to insert a note in the text itself, rather than in the margin.
\begin{Verbatim}
	\todo[inline]{Check the experimental setup before submission.}
\end{Verbatim}
Output:
\todo[inline]{Check the experimental setup before submission.}

\vspace{0.5cm}\noindent You can also specify the author and color of the note:
\begin{Verbatim}
	\todo[inline,author=Jacopo,color=green!20]{Update the results.}
\end{Verbatim}
Output:
\todo[inline,author=Jacopo,color=green!20]{Update the results.}

\section*{List of todos}

To generate a list of all todos in the document, use the \verb|\listoftodos[<Custom Title>]| command at the desired location.It can also be added to the Table of Contents using \verb|\todototoc| in the preamble. Example below:

{
	\hypersetup{hidelinks}
	\listoftodos[My Todos]
}

\end{document}
